\scanpage{54}
\mspage{10}
of weighted graphs with markings, the operation \deleted{\unsure{is made explicit}}
is reflected in a corresponding operation at the level of $\vec R$, by the ``erasure'' of a subset $\vec R'$ stable under $\mathbb D_\infty$ i.e. under $\sigma_0,\sigma_1$.

2) Erasure of cutting circles.

We are now given a \deleted{set of connected components of} 1-submanifold
\inserted{(open and)} \inserted{closed} $\Sigma'$ of $\Sigma\setminus\partial X$,
corresponding to a set \deleted{\unsure{A}} of connected components of $\Sigma'$ i.e. a subset \inserted{$A'_0$} of $\pi_0(\Sigma')=A_0$
(ordinary edges of the graph) --- and we replace $\Sigma$ by $\Sigma\setminus\Sigma'$.
At the level of weighted graphs, this translates into an operation of \uline{contraction to a point} of the corresponding edges, hence formation of a new graph, and of a \uline{generization morphism} (as opposed to ``specialization'') from the old to the new.
It is appropriate here to make explicit how the weighting follows\ldots\
At the level of the \inserted{vertices} (markings), it is understood that the marked vertices lying on $\Sigma'$ are erased at the same time --- thus, on the weighted graph with markings, we remove the set of markings which are over the edges $\in A'_0$.
\mspage{53}
